\chapteropen[\href{https://en.wikipedia.org/wiki/Automaton}{Automaton in the Swiss museum Centre International de la M\'ecanique d'Art, writing a letter.}]{automata/automaton.jpg}{Automata} %https://commons.wikimedia.org/wiki/File:CIMA_mg_8332.jpg
% License of photo: This work is free software; you can redistribute it or modify it under the terms of the CeCILL. The terms of the CeCILL license are available at www.cecill.info.

% Use as much of http://cstheory.stackexchange.com/a/14818/4731 as possible


\section{Finite state automata}

\subsection{Definition}


\subsection{Kleene's Theorem}

\subsection{Pumping}
Pumping Lemma

\subsection{Minimization}

Motivation

\begin{definition}
Let $\lang{L}$ be a language over an alphabet~$\Sigma$.
For a pair of strings $\sigma,\hat{\sigma}\in\kleenestar{\Sigma}$,
a \definend{distinguishing extension} is a 
string $\tau\in\kleenestar{\Sigma}$ 
such that exactly one of the two strings $\sigma\concat\tau$ 
and $\hat{\sigma}\concat\tau$ is an element of~$\lang{L}$. 
\end{definition}

Examples, including example where the two do not have a distinguishing 
extension.

\begin{definition}
Two strings $\sigma,\hat{\sigma}\in\kleenestar{\Sigma}$ are 
definend{$\lang{L}$-related}, written $\sigma\sim_{\lang{L}}\hat{\sigma}$,
if they have no distinguishing extension.  
\end{definition}

\begin{lemma}
For any language~$\lang{L}$,
the binary relation $\sim_{\lang{L}}$ between strings
is an equivalence relation.   
\end{lemma}

\begin{proof}
Sketch, including reminder of what is an equivalence.  
\end{proof}

Examples
Reminder of equivalence classes.

\begin{theorem}[Myhill-Nerode]
A language~$\lang{L}$ is regular if and only if the relation
$\sim_{\lang{L}}$ has only a finite number of equivalence classes.
\end{theorem}

Moreover, the proof will show that there is a correspondence between the classes
and the states of a finite state automata accepting that lanuguage.

\begin{proof}
Suppose that $\lang{L}$ is a regular, over the alphabet~$\Sigma$.
Then there is a finite state automata $\mathcal{A}$ accepting~$\lang{L}$.
Partition the set of all strings~$\kleenestar{\Sigma}$ with:
two strings $\sigma,\hat{\sigma}$ are in the same part if 
when~$\mathcal{A}$ is given them as input then they take 
the machine to the same state.
Observe that if two strings are in the same part then they are 
$\lang{L}$-related.
Thus, the parts are subsets of the equivalence classes of 
$\sim_{\lang{L}}$.
Since the number of parts is less than or equal to the number of states
in $\mathcal{A}$, the number of equivalence classes is finite.

If L is a regular language, then by definition there is a DFA A that recognizes it, with only finitely many states. If there are n states, then partition the set of all finite strings into n subsets, where subset Si is the set of strings that, when given as input to automaton A, cause it to end in state i. For every two strings x and y that belong to the same state, and for every choice of a third string z, automaton A reaches the same state on input xz as it reaches on input yz, and therefore must either accept both of the inputs xz and yz or reject both of them. Therefore, no string z can be a distinguishing extension for x and y, so they must be related by RL. Thus, Si is a subset of an equivalence class of RL. Combining this fact with the fact that every member of one of these equivalence classes belongs to one of the sets Si, this gives a many-to-one relation from states of A to equivalence classes, implying that the number of equivalence classes is finite and at most n.


\end{proof}

\section{Pushdown automata}

\subsection{Definition}

\extra{Neural nets}